% Generated once by scripts/paper_md2tex.py from docs/paper/draft.md; edit this file from now on.
\section{Related work in full}\label{app:H}
\noindent Record section 2.
\subsection*{2. Related work and boundaries}\addcontentsline{toc}{subsection}{2. Related work and boundaries}

\href{https://arxiv.org/abs/2106.09685}{LoRA} supplies the starting parameterization: a frozen
base projection plus a trained low-rank branch. Our question concerns context-conditioned
computation inside that branch. It does not depend on presenting recurrence itself as a
new addition to low-rank adaptation. The dated literature audit
is the scope record for the comparisons below; we retain its access limitations rather
than turning an incomplete search into a priority claim.

Associative fast weights are the closest antecedent to the memory access pattern.
\href{https://arxiv.org/html/1610.06258v3}{Ba et al.} already use an evolving inner state to
read fixed fast weights repeatedly, with sustained input during the inner loop. Their
construction makes clear that repeated access to a fixed matrix predates this adapter
study. \href{https://proceedings.mlr.press/v139/schlag21a.html}{Schlag et al.} connect linear
attention to fast-weight programming and motivate delta-style updates that correct
stored key-value associations. We use that lineage to separate the slow parameters
that learn a representation from the online matrix that stores the current context.

\href{https://arxiv.org/html/2605.12357v2}{\ensuremath{\delta}-mem} combines associative online state with
low-rank corrections to a frozen language model. Its old-state read provides a useful
comparison with our separation between writing the context and spending several reader
updates on the resulting snapshot. The relevant distinction is the read schedule and
query evolution, not a claim that associative memory or context-dependent low-rank
corrections are new. Matching these schedules also requires specifying the writer’s
decay/\allowbreak{}update order, because nominally similar delta rules can retain different states.

Looped models investigate where repeated computation should occur.
\href{https://arxiv.org/abs/2608.18230}{MixerLoop} repeats a delta-network mixer while applying
its dense feed-forward component once; the literature audit distinguishes its repeated
mixer updates from the fixed-memory reader considered here. \href{https://arxiv.org/abs/2503.07903}{MemReasoner}
is a close memory-and-composition neighbor, with an end-to-end trained memory-augmented
language model. These comparisons motivate describing the exact object being iterated:
a low-dimensional query, a whole hidden state, or an update that also rewrites memory.
They do not justify transferring a result between those objects without an equivalence
argument.

\href{https://arxiv.org/abs/1909.01377}{Deep equilibrium models} motivate another distinct
question: finding and differentiating through a fixed point. Our finite-depth reader
need not be a contraction to be useful, and a fixed-point error certificate would not
by itself establish correct task predictions. RecLoRA remains a related-work item whose
full-text architectural comparison is \textbf{TODO}: the committed audit could not verify its
page-level claims. We consequently make no claim about absent nonlinearities, routing,
or token mixing in that work, and do not use its unverified publication status in the
argument. The contribution sought here is a reproducible composition-and-budget study
with explicit writer controls and conditional analysis.

\subsubsection*{Latent recurrence, sharing, and stopping}

The \href{https://arxiv.org/html/2507.06203v2}{latent-reasoning survey} distinguishes
vertical activation recurrence at fixed token/\allowbreak{}time from evolution across sequence
history. Our reader uses the former axis while the writer updates memory across
tokens. \href{https://arxiv.org/html/2502.05171v2}{Huginn} supplies precedents for shared
Transformer recurrence, reinjected input, heavy-tailed training-depth sampling,
and heuristic prediction-KL exits. \href{https://arxiv.org/html/2412.06769v4}{Coconut}
instead adds latent sequence positions by reusing hidden states as input embeddings;
its default fixed latent length should not be confused with guaranteed adaptive
halting. These are different placements of latent computation, with different caches
and training objectives.

\href{https://arxiv.org/pdf/2510.25741v6}{Ouro} learns exit allocation for shared Transformer
loops and reports degradation beyond trained depth. That is a precedent for studying
nonmonotone depth, not a diagnosis of this reader’s failure mechanism.
\href{https://arxiv.org/html/2410.20672v3}{Relaxed Recursive Transformers} use depth-specific
LoRA to relax full-block sharing and train all model parameters; their oracle-exit
batching simulations do not establish this adapter’s measured efficiency.
\href{https://arxiv.org/html/2507.10524v3}{Mixture-of-Recursions} adds learned token-level
depth routing, including a variant that reuses first-recursion KV for later queries.
Consequently fixed context alone is not a unique distinction. Our specified compact
delta state, low-dimensional reader, and separate write schedule define the comparison.

Degradation from extra iterations is a known failure of recurrent depth, separate from this
adapter. \href{https://arxiv.org/abs/2202.05826}{Bansal et al.} report that recurrent networks can
lose accuracy when iterated beyond their training depth (“overthinking”), and reduce it with a
recall path that reinjects the input together with progressive training. The reader’s
reinjected u is analogous to a recall path, and no progressive or stay objective is trained
here. \href{https://arxiv.org/abs/2511.08577}{Think-at-Hard} reports that extra loop iterations can
turn correct tokens into errors in looped language models, and
\href{https://arxiv.org/abs/2604.07822}{Kohli et al.} report that recurrent-depth transformers can
compose knowledge, need more recurrence for deeper compositions, and overthink under excessive
recurrence. We therefore read the loss of short-hop accuracy at larger L (Sections 5.1 and 7)
as an instance of a known failure in a new setting, not as a defect specific to this reader,
and we do not test the remedies those works propose.  % S7 JS1

\href{https://arxiv.org/html/2605.26733v1}{STARS} already combines sampled depth with
Jacobian-vector regularization of looped dynamics. Its practical sampled-state proxy
is not a certified spectral radius or global contraction constant.
\href{https://arxiv.org/html/2608.08113v1}{ReLIT} combines a frozen-backbone anchor with
recursive full-sequence answer/\allowbreak{}scratchpad updates and ACT halting. These precedents
narrow any claim about reinjection, frozen-backbone latent recurrence, or stabilization.
They motivate the JS protocol; its separate adapter results are
reported in Section 7.

\subsubsection*{Multi-hop composition, regression memories, and retrieval in weights}

The works in this subsection were added after two external literature reviews
(rv-fable and rv-astra)
opened their primary pages. We cite only what their abstracts and the review notes support,
re-read no full text for this manuscript, and draw no priority claim from them.

Latent multi-hop composition in language models has a substantial literature.
\href{https://arxiv.org/abs/2402.16837}{Yang et al. (2024)} look for and intervene on the bridge
entity of two-hop queries, and \href{https://arxiv.org/abs/2411.16679}{Yang et al. (2025)} remove
shortcuts and find relation-dependent composition.
\href{https://arxiv.org/abs/2411.16353}{Balesni et al.} report both failures and successes of
latent two-hop reasoning across synthetic and natural settings.
\href{https://arxiv.org/abs/2406.12775}{Biran et al.} report that the bridge is resolved early and
the second hop starts too late in the layer stack,
\href{https://arxiv.org/abs/2502.10835}{Yu et al.} let early layers read later-layer states, and
\href{https://arxiv.org/abs/2601.03542}{Liu et al.} report that later-hop answers can be decodable
before the bridge. \href{https://arxiv.org/abs/2502.13913}{Guo et al.} find in-context two-hop
reasoning with distractors near guessing before fine-tuning. For out-of-distribution
composition, \href{https://arxiv.org/abs/2608.07261}{Zhang et al.} report failure when upper layers
map facts to outputs instead of reasoning, and that recurrent-style training that reuses
circuitry helped there; \href{https://arxiv.org/abs/2509.24653}{Lin et al.} report that identity
supervision on bridges enables unseen compositions, and
\href{https://arxiv.org/abs/2405.15071}{Wang et al.} that composition does not generalize out of
distribution without grokking. Our setting differs: facts reach the reader only through a
fixed-size adapter memory, and hops are spent in a reader loop rather than in layer depth.
These studies therefore bound what we may claim without predicting our results. In particular,
Zhang et al. make it unsafe to read B8’s in-distribution advantage over the shared reader
(Section 6; B8 is also not parameter-matched) as evidence for the unshared design, because that advantage is measured on trained
hops and names; and the paper reports only a descriptive scope row on unseen-pool names of its
closed vocabulary (Section 9), not a test of composition over new names or codes.
% B3 B6 UN2 UN3

Associative memory can also be viewed as regression at test time.
\href{https://arxiv.org/abs/2501.12352}{Wang et al.} treat a delta-rule write as one gradient step
on a recall loss, with least squares as the exact minimizer.
\href{https://arxiv.org/abs/2406.06484}{Yang et al.} parallelize delta-rule linear transformers over
sequence length, and \href{https://arxiv.org/abs/2501.00663}{Titans},
\href{https://arxiv.org/abs/2505.23735}{ATLAS}, and \href{https://arxiv.org/abs/2407.04620}{TTT layers}
train the memory itself at test time, ATLAS over a window of past tokens rather than only the
latest one. Our writer is the one-step delta update, which need not reach the regression
solution when keys are correlated (cf. Proposition 11). Whether that gap contributes to the
deployed-read distortion in Section 8.3 is untested, because the audit identifies no unique
cause, and we study none of these alternative writers.  % LL7 LL9

Slot and lookup memories give each stored item its own address.
\href{https://arxiv.org/abs/1503.08895}{End-to-end memory networks} learn several memory hops over
stored items, and \href{https://aclanthology.org/D16-1147/}{key-value memory networks} separate the
addressing representation from the returned one, so multi-hop reading of a memory and
key-value separation are long established and we claim neither.
\href{https://arxiv.org/abs/2008.02217}{Hopfield networks} retrieve stored patterns in a single
softmax step, \href{https://arxiv.org/abs/2412.09764}{memory layers} scale trainable key-value
lookup inside language models, and \href{https://arxiv.org/abs/2405.06394}{Memory Mosaics} compose
networks of associative memories. Their state grows with the number of stored items, whereas
the reader here uses a fixed-size delta-rule state in which facts superpose. No slot-memory
baseline is run, and no result here bears on how a slot memory would behave.

Facts can also live in, or be written into, parameters.
\href{https://arxiv.org/abs/2012.14913}{Geva et al.} interpret feed-forward layers as key-value
memories; \href{https://arxiv.org/abs/2202.05262}{ROME} edits a factual association with a
rank-one update and \href{https://arxiv.org/abs/2210.07229}{MEMIT} scales such edits to many facts,
a linear associative write into persistent weights rather than into a resettable per-episode
state. \href{https://arxiv.org/abs/2403.11901}{Larimar} attaches an episodic memory with one-shot
writes to a language model; \href{https://arxiv.org/abs/2506.06266}{Cartridges} and
\href{https://arxiv.org/abs/2411.09289}{StreamAdapter} compress a context into a trained
key-value cache and a parameter update, respectively, Cartridges by training on each context,
whereas the writers studied here are single-pass delta updates. Three 2026 reports on
memories for frozen language models locate the bottleneck in reading or extraction rather
than in storage (\href{https://arxiv.org/abs/2608.17050}{Frozen Memory Is Not Enough},
\href{https://arxiv.org/abs/2601.15324}{Prometheus Mind}, and a workshop paper on
\href{https://neurips.cc/virtual/2026/170472}{a hidden-state memory layer}). That is compatible
with, but does not confirm, the read-distortion measurements of Section 8.3, whose cause the
audit leaves unidentified.  % LL7 LL9

Two further external reviews (rv-fable and rv-astra, 2026-10-29) added the works
below. We cite their abstract pages as the reviewers recorded them, did not re-open the pages for this manuscript, and claim no priority from any of them.
We study none of the writers, objectives or architectures named here; our writers are the one-step delta writer, its DPFP-feature variant and the oracle writers.
\href{https://arxiv.org/abs/2601.15313}{Zahn et al.} report that semantic keys overlap and recall falls near chance after a few related facts, while identifier keys do
not, a setting close to the key-overlap questions of Section 8.3. Several works read the delta rule as online regression or filtering:
\href{https://arxiv.org/abs/2605.11196}{Pandey and Singh} use online ridge by Sherman-Morrison as approximate key orthogonalization;
\href{https://arxiv.org/abs/2511.21016}{Gated KalmaNet} shows a full-covariance Kalman update to be online ridge regression with fading memory;
\href{https://arxiv.org/abs/2602.10743}{Kalman Linear Attention} filters in information form, scan-parallel;
\href{https://arxiv.org/abs/2609.07816}{Kalman Delta Networks} use a per-channel Kalman gain whose isotropic case is the delta rule;
\href{https://arxiv.org/abs/2604.21100}{Preconditioned DeltaNet} reads the delta rule as a first-order regression step and adds preconditioning;
\href{https://arxiv.org/abs/2407.14207}{Longhorn} derives a state-space update as closed-form online recall;
\href{https://arxiv.org/abs/2506.05233}{MesaNet} optimizes a local regression at test time; and
\href{https://arxiv.org/abs/2606.08804}{Q-Delta} studies keys transformed by subsequent updates. Other works change the objective, the retention or the update:
\href{https://arxiv.org/abs/2504.13173}{It’s All Connected} (attentional-bias objectives and retention gates),
\href{https://arxiv.org/abs/2512.24695}{Nested Learning} (a self-modifying model with continuum memory),
\href{https://arxiv.org/abs/2512.23675}{End-to-End Test-Time Training} (a meta-learned initialization),
\href{https://arxiv.org/abs/2505.23884}{Test-Time Training Done Right} (whole-chunk updates),
\href{https://arxiv.org/abs/2404.07143}{Infini-attention} (compressive memory with a delta-style update),
\href{https://arxiv.org/abs/2605.22791}{Gated DeltaNet-2} (separate erase and write gates),
\href{https://arxiv.org/abs/2610.02816}{Gated Slot Attention-2} (two-sided association correction through slots),
\href{https://arxiv.org/abs/2502.10297}{DeltaProduct} and \href{https://arxiv.org/abs/2607.07386}{Sparse Delta Memory} (richer updates or larger states).
\href{https://arxiv.org/abs/2402.13449}{CAMELoT} attaches a training-free associative memory to a frozen language model, as a neighbor of the frozen-backbone
setting here. These are different writers from those tested, and the learned-writer stop of Section 8.5 applies only to the writers it tested.

For looped and latent reasoning, \href{https://arxiv.org/abs/2607.00341}{DiscoLoop} reports that a bridge can be decodable yet misaligned for the next loop and mixes
decoded embeddings with continuous states; \href{https://arxiv.org/abs/2606.31779}{Fan, Svete and Lee} (called LOTUS in rv-astra) supervise looped transformers per position;
\href{https://arxiv.org/abs/2603.21676}{Chen} trains with final-only supervision, a near-identity gate and depth extrapolation, where intermediate supervision hurt;
\href{https://arxiv.org/abs/2606.29983}{Kuo et al.} report that a randomized loop count stabilizes extrapolation; and
\href{https://arxiv.org/abs/2502.21074}{CODI} and \href{https://arxiv.org/abs/2610.10623}{Recurrent Self-Improvement} distill explicit reasoning into latent reasoning or
use deeper teachers. \href{https://arxiv.org/abs/2402.09268}{Sanford et al.} show that logarithmic depth suffices for k-hop problems, so one hop per read is a design
choice of this adapter and not a bound. \href{https://arxiv.org/abs/2602.00449}{Liang and Pan} warn that decodable intermediate states do not establish a complete serial
algorithm, which is the boundary the per-step readouts of Section 8.5 respect: they are decoding only. \href{https://arxiv.org/abs/2607.19390}{The orthogonalized read as a training scaffold}  % LP5
is reported as a removable training aid and not an inference rescue; whether sharing acts as such a scaffold for B8 is untested here. None of these
training objectives, stopping rules or distillation methods is used in this paper.

\subsubsection*{Verbalizable directions and the readout boundary}

\href{https://transformer-circuits.pub/2026/workspace/index.html}{Verbalizable Representations Form a Global Workspace in Language Models}
defines J-lens using averaged transport from intermediate residuals toward final-layer
representations and vocabulary readout. Its sparse J-space construction is generally
a union of cones from an overcomplete frame, not the adapter’s fixed col(B) or a single
rank-r subspace. It motivates asking whether intermediate reader states expose the
expected entities, while separating readability from causal use. The downstream
transport Jacobian and the reader’s recurrent transition Jacobian are different
objects. The expanded primary-source audit
records versions, access limits, and these distinctions; it supports no new priority claim.

